\begin{table}[t]
\centering
\caption{Summary Statistics: FEMA IHP Disasters, 2010--2025}
\label{tab:summary}
\begin{tabular}{lrrrr}
\toprule
 & Mean & SD & P25 & P75 \\
\midrule
\multicolumn{5}{l}{\textit{Panel A: Disaster level} ($N = 364$)} \\[3pt]
Declaration lag (days) & 35.4 & 33.3 & 7.0 & 63.0 \\
Concurrent disasters (30d window) & 25.7 & 25.0 & 5.0 & 57.0 \\
Recent declarations (60d) & 4.0 & 6.6 & 1.0 & 5.0 \\
Total registrations & 25,223 & 102,129 & 755.2 & 8,776 \\
IHP per registrant (\$) & 3,439 & 2,952 & 1,602 & 4,729 \\
Approval rate & 0.5 & 0.2 & 0.3 & 0.7 \\
Avg.\ FEMA-inspected damage (\$) & 5,021 & 9,409 & 1,365 & 5,338 \\
Counties affected & 19.1 & 28.1 & 4.0 & 22.0 \\
\addlinespace[6pt]
\multicolumn{5}{l}{\textit{Panel B: City--disaster level} ($N = 57,175$)} \\[3pt]
IHP per registrant (\$) & 2,321 & 3,670 & 0.0 & 3,091 \\
Approval rate & 0.455 & 0.380 & 0.000 & 0.800 \\
FEMA-inspected damage (\$) & 2,097 & 6,254 & 0.0 & 2,067 \\
Valid registrations & 160.6 & 1,697 & 1.0 & 24.0 \\
\bottomrule
\end{tabular}
\begin{tablenotes}[flushleft]
\small
\item \textit{Notes:} Data from OpenFEMA Individual Assistance program, 2010--2025.
Declaration lag is the number of days between incident onset and presidential disaster
declaration. IHP = Individual and Households Program.
\end{tablenotes}
\end{table}
